\clearpage\thispagestyle{otFront}\phantomsection\label{references}\pdfbookmark[0]{Sources et édition}{bm-references}\otPageTitre{Sources et édition}{Références officielles · France}\fontsize{11}{15}\selectfont
\par\medskip\textbf{1. Programme applicable}\par Bulletin officiel du 30 juillet 2020 : programme du cycle 4, encore applicable en Quatrième en 2026-2027.\par\href{https://www.education.gouv.fr/bo/20/Hebdo31/MENE2018714A.htm}{\textcolor{otGreenDark}{Consulter la source officielle}}\par
\par\medskip\textbf{2. Attendus de fin de Quatrième}\par Éduscol : document complet de 12 pages, connaissances et exemples de réussite, y compris programmation de niveaux 1 et 2.\par\href{https://eduscol.education.gouv.fr/sites/default/files/document/16-maths-4e-attendus-eduscol1114746pdf-74682.pdf}{\textcolor{otGreenDark}{Consulter la source officielle}}\par
\par\medskip\textbf{3. Repères annuels du cycle 4}\par Éduscol : progression par niveau, limites du formalisme, géométrie et programmation.\par\href{https://eduscol.education.gouv.fr/sites/default/files/document/26-maths-c4-reperes-eduscol1114756pdf-74694.pdf}{\textcolor{otGreenDark}{Consulter la source officielle}}\par
\par\medskip\textbf{4. Calendrier du nouveau programme}\par BO du 5 mars 2026, article 3 : nouveau programme applicable en Quatrième à la rentrée 2027-2028.\par\href{https://www.education.gouv.fr/bo/2026/Hebdo10/MENE2602912A}{\textcolor{otGreenDark}{Consulter la source officielle}}\par
\par\medskip\textbf{5. Validité des attendus en 2026-2027}\par Éduscol, repères et attendus annuels : calendrier de leur remplacement progressif.\par\href{https://eduscol.education.gouv.fr/6910/reperes-annuels-de-progression-et-attendus-de-fin-d-annee-du-cp-la-troisieme}{\textcolor{otGreenDark}{Consulter la source officielle}}\par
\par\bigskip Sources consultées le 22 octobre 2026. Les exercices et bilans
sont originaux, sans reproduction des exemples officiels.\par\medskip
Périmètre : Thalès dans les triangles emboîtés, cosinus seul en trigonométrie,
translations sans calcul vectoriel formel, dépendance de grandeurs sans
notations de fonctions, probabilités à une épreuve, programmation de niveaux
1 et 2. Les défis approfondissent ces outils sans ajouter d’exigence officielle.
Le programme nouveau de 2027-2028 en Quatrième demandera une édition adaptée.

\clearpage
